\pdfoutput=1
\documentclass[11pt]{amsart}
\usepackage{amsmath,amssymb,amsthm,microtype}
\usepackage[margin=1.2in]{geometry}
\usepackage{xcolor}
\usepackage[colorlinks=true,linkcolor=blue,citecolor=blue,urlcolor=blue]{hyperref}
\hypersetup{pdfauthor={Bernd Johannes Wuebben},pdftitle={Integral generators for curve counting on threefolds}}
\setlength{\emergencystretch}{2em}
\newtheorem{theorem}{Theorem}[section]
\newtheorem{proposition}[theorem]{Proposition}
\newtheorem{lemma}[theorem]{Lemma}
\newtheorem{corollary}[theorem]{Corollary}
\theoremstyle{remark}\newtheorem{example}[theorem]{Example}
\newcommand{\ZZ}{\mathbb Z}
\newcommand{\QQ}{\mathbb Q}
\newcommand{\CC}{\mathbb C}
\newcommand{\PP}{\mathbb P}
\newcommand{\cO}{\mathcal O}
\newcommand{\tot}{\operatorname{Tot}}
\newcommand{\Sym}{\operatorname{Sym}}
\title[Integral generators for curve counting]{Integral generators for curve counting on threefolds}
\author{Bernd Johannes Wuebben}
\subjclass[2020]{14N35 (primary); 14C25, 14J32, 53D45 (secondary)}
\keywords{Calabi--Yau threefolds, universal curve counting, local curves, integral generators, conifold transitions}
\date{November 24, 2026\\This version: November 24, 2026}
\begin{document}
\begin{abstract}
Pardon's integral curve-counting algebra has geometric generators
whose local models may depend on the covering multiplicity. In first
Chern number and virtual dimension zero we give a single model in
each genus for all multiplicities. On a genus-$g$ curve take two line
bundles of degree $g-1$ with canonical determinant and act on them
with opposite weights. The equivariant classes of the multiple zero
sections freely generate the integral algebra. The comparison with an
isolated holomorphic multiple is triangular in multiplicity, with
leading coefficient one given by the complex orientation of the embedded
curve. The resulting factors give a unique product with integral
exponents for effective cycles of Chern number zero on every compact
Kähler threefold whose first Chern class is nonnegative on curves.
\end{abstract}
\maketitle
\section{The integral local factors}\label{sec:intro}

The universal count of effective curves on a compact threefold records
the entire space of cycles in each integral homology class. Even when
these spaces carry integral classes, the factors used to express their
generating series can have fractional exponents. On a smooth projective
Calabi--Yau threefold, for example, a genus-ten factor occurs with
exponent in $\mathbb Z+2/3$ in three times a primitive class
\cite[Theorem~1.1]{WueCounter}, \cite[Theorem~1.2]{WueTriples}. This factor comes from scaling the two
normal directions to a curve by the same weight. Which local factors
recover the integral structure?

Pardon constructed an algebra of effective one-cycles in families of
complex threefolds. In its semi-Fano part every irreducible component
has nonnegative first Chern number. Write $H_{\ZZ}$ for the summand
in first Chern number and virtual dimension zero. Multiplication is
external product on disjoint targets. Pardon proved that integral
geometric elements obtained by isolating multiples of embedded curves,
one for each genus and multiplicity, freely generate this algebra.
The normal neighborhood used to isolate a multiple may depend on its
multiplicity. A uniform construction by scalar dilation gives free
generators over $\QQ$ \cite[Theorem~1.2]{Pardon}. We identify a uniform
equivariant construction which gives free generators over $\ZZ$.

Let $C$ be a smooth compact curve of genus $g$. Choose a line bundle
$L$ of degree $g-1$ and give the two summands of
$L\oplus(K_C\otimes L^{-1})$ weights $+1$ and $-1$. We call this
the balanced antidiagonal model. Write $a_{g,m}$ for the normalized
equivariant class supported on the $m$-fold zero section and
$A_g(q)=1+\sum_{m\ge1}a_{g,m}q^m$. This class is independent of the
choices of $C$ and $L$. The normalization divides a supported Thom
class by a product of weights, so its construction gives an element of
$H=H_{\ZZ}\otimes\QQ$. For equal weights write $x_{g,m}$ for the scalar
class. The two classes agree in multiplicity one.

\begin{theorem}\label{thm:basis}
The map to the rationalization embeds $H_{\ZZ}$ as the polynomial ring
$\ZZ[a_{g,m}:g\ge0,\ m\ge1]$. For each $(g,m)$ one can choose an
integral geometric local curve element $y_{g,m}$ of this type so that
\[
 y_{g,m}-a_{g,m}\ \in\ \ZZ[a_{h,j}:h\ge0,\ 1\le j<m].
\]
\end{theorem}

There is a geometric obstruction to comparing the two sorts of
generators directly in the normal bundle. For $g\ge2$ choose $L$ outside
the theta divisor. Both normal summands are acyclic, so the embedded
zero section is regular and isolated. Yet $h^0(C,L^2)\ge g-1$.
A nonzero section $s$ cuts out the compact double cycles $w^2=ts$ in
the first summand. They converge to the doubled zero section as $t$
tends to zero. Isolation of the embedded curve says little about
isolation of its doubles.

For a prescribed multiplicity $m$ we start in a holomorphic
neighborhood in which all multiples through $m$ are isolated. We place a tube of this neighborhood in a closed symplectic six-torus.
Bounded-area compactness keeps the cycles in the tube away from its
boundary during a small deformation. Excision therefore lets us apply
the compact cluster comparison to its insertion, giving a product of contributions of elementary curves, the regular local
models of Ionel and Parker \cite{IP}. A genus-$h$ elementary curve
contributes the balanced series $A_h$ by the geometric comparison of
\cite[Proposition~3.1]{WueProduct}. In degree one, testing by genus
separates the multiplicity-one classes: their fixed-genus series have
distinct leading powers. The embedded curve has its complex orientation,
so its genus occurs with coefficient $+1$ and every other genus with
coefficient zero. The remaining factors have core
class a proper multiple of the original class. In the coefficient of
$q^m$ these corrections therefore involve only generators of
multiplicity less than $m$. This gives the integral triangular
comparison in Theorem~\ref{thm:basis}.

For a smooth compact Kähler threefold $X$, write $\Gamma_0(X)$ for zero
and the effective integral classes of first Chern number zero. Let
$U_{X,\beta}$ be the characteristic function of the entire space of
effective cycles of class $\beta$, and take $U_{X,0}=1$.
\begin{theorem}\label{thm:compact}
If $c_1(TX)$ is nonnegative on every compact irreducible curve, then
\begin{equation}\label{eq:compact}
 \sum_{\beta\in\Gamma_0(X)}U_{X,\beta}Q^\beta
 =\prod_{\substack{\beta\in\Gamma_0(X)\\ \beta\ne0}}\prod_{h\ge0}
           A_h(Q^\beta)^{n_{\beta,h}}
 \quad\text{in }H_{\ZZ}[[\Gamma_0(X)]] .
\end{equation}
The exponents are unique integers and only finitely many occur for each
$\beta$. The completion is by Kähler area, and the integers are independent
of the Kähler form. When $X$ is smooth projective with trivial canonical
bundle and $H^1(X,\cO_X)=H^2(X,\cO_X)=0$, the exponent $n_{\beta,h}$
is the signed multiplicity of elementary clusters of class $\beta$ and
genus $h$ in the bounded-energy cluster decomposition.
\end{theorem}

Two identities for cycles explain why the basis theorem yields a
product. Addition records every decomposition of a cycle, and division
by an integer records every integral root of its class. After the
lower-area factors have been removed, the first remaining coefficient
is primitive for addition and vanishes under every proper division.
It is therefore an integral combination of multiplicity-one classes.
Their coefficients give the exponents at that homology label.
The product retains the full integral homology label, including
torsion. Its local factors also have a numerical interpretation.
Bryan and Pandharipande computed the local Gromov--Witten residues
\cite{BP}, and Ionel and Parker identified the antidiagonal residue of an
elementary curve \cite[Proposition~3.3]{IP}. Doan, Ionel and Walpuski
proved the requisite finiteness for the symplectic cluster expansion
\cite{DIW}. The balanced universal classes and a conditional numerical
factorization were constructed in \cite{WueBalanced}. The geometric
comparison of \cite{WueProduct} established the projective cluster
formula. Here its use on a tube gives the integral presentation and
the compact Kähler factorization. In the compact example above, the
primitive genus-three count is $-14$. Changing back to scalar factors
subtracts $(4720/3)(-14)$ from the genus-ten exponent, giving the stated
fractional part \cite{WueCounter,WueTriples}.

\subsection*{Acknowledgements}
I thank John Pardon for pointing out the difficulty of comparing his
scalar-equivariant local curve classes with the elementary curves of
Ionel and Parker.

\section{Comparing local curve classes}\label{sec:comparison}

\subsection{Two normal weights}
A degree-$m$ cycle in the total space of $L_+\oplus L_-$, with
$L_+=L$ and $L_-=K_C L^{-1}$, meets a fiber of the projection to $C$
in a multiset of length $m$. No fiber contains a compact curve.
The elementary symmetric functions of either coordinate in this multiset
are sections of $L_\pm^j$, $1\le j\le m$. For instance, on the
normalization of a component the characteristic polynomial of
multiplication by the fiber coordinate on its finite flat direct image
extends the symmetric functions across branch values. Multiply the
polynomials with their cycle multiplicities. Evaluation at more than
$\max(0,m(g-1))$ distinct points therefore gives a tuple with common
zero precisely $m[C]$. Its weights are the integers $\pm j$.

For clarity, a tuple of $k$ nonzero weights $r_1,\ldots,r_k$ on the
cycle space gives, over the finite Borel base $\PP^k$, a section of the
pullback of $\bigoplus_i\cO(r_i)$. Pull back its relative Thom class and
divide by $\prod_i r_i$. The support on the relative cycle space is
compact. Adding an equivariant test after passage to the next projective
space gives the hyperplane pushforward of this class, since the normalized
Euler class of the added line is the hyperplane class. Concatenating two
tuples proves independence of the tests. Varying the curve and line of
degree $g-1$ gives the same universal class $a_{g,m}$
\cite[\S\S2--3]{WueBalanced}. The division in this prescription explains why
integrality of its answer requires proof.

The first two weights can already distinguish the prescriptions. Over a
curve of genus two the scalar and balanced doubles satisfy
$a_{2,2}-x_{2,2}=-2x_{4,1}$. In the holomorphic double-cover
calculation two smooth branches of genus four move away from the
zero section. Their contribution fixes the sign in the comparison
\cite[\S8]{WueProduct}.

There are two operations on the resulting coefficients. Pullback by
addition of effective cycles is the coproduct $\Delta$. Pullback by
multiplication of a cycle by $l$ is denoted $\rho_l$. With $u_0=1$, the
scalar and the antidiagonal classes of a fixed genus both satisfy
\begin{equation}\label{eq:localops}
 \Delta u_d=\sum_{i+j=d}u_i\otimes u_j,\qquad
 \rho_l u_d=\begin{cases}u_{d/l}&l\mid d,\\0&l\nmid d.\end{cases}
\end{equation}
These identities include the diagonal of Borel bases in the coproduct.
In each shifted term the two relative cohomological degrees are $d$
and $-d$, with $d\ne0$. In Chern number zero the factor of positive
relative cohomological degree has negative virtual dimension and vanishes
\cite[Lemma~4.13 and Theorem~6.13]{Pardon},
\cite[Proposition~3.1]{WueBalanced}. The augmentation is restriction to the
empty cycle.

The laws \eqref{eq:localops} alone leave a choice of integral
coordinates. For instance one may replace $A_g(q)$ by
$A_g(q)A_h(q^2)$ for $h\ne g$ and invert the change integrally.
The unit of a positive elementary curve specifies the particular
coordinates used in Theorem~\ref{thm:compact}.

\subsection{Elementary curves and multiplicity one}
We use the comparison $\Phi:H\to S$ of \cite[\S2]{WueProduct}, where $S$
is the rational cycle algebra for smooth families of almost-complex
six-manifolds. Over a point the characteristic function of a compact
open set of cycles gives an insertion. Open extension in the target,
external product for disjoint targets, and the two endpoint restrictions
of a compactly supported class over an interval are the operations we
will need. The geometric topology records the support and the integration
current of the cycle. In integrable families it gives the topology of
compact analytic cycles.

We recall three results from \cite{WueProduct} and their roles in the
comparison. First, on the elementary genus-$h$ curve of Ionel and Parker
the unit at its $j$-fold cycle is $\Phi(a_{h,j})$
\cite[Proposition~3.1]{WueProduct}. On the sum of two copies of a spin line,
real rotations commute with the Lee deformation to a super-rigid normal
operator; in complex eigenlines those rotations have weights $+1,-1$.
The normalized Thom class extends on the finite Borel family, restricting
at the elementary end to the unit of the isolated cycle. Second,
a bounded-energy partition into clusters gives a \emph{product} of these
units and their inverses \cite[\S4, Proposition~4.1]{WueProduct}. It is
essential here to partition disconnected cycles as well as connected
maps. We explain its local use in \S\ref{sec:basis}.

Third, the classes $\Phi(x_{h,1})$ are linearly independent over $\QQ$
\cite[Theorem~5.1]{WueProduct}. The argument tests a proposed finite
relation at its least genus $h$. The image-cycle map on connected stable
maps of that genus is proper over compact sets of cycles. A class supported
on an embedded curve of larger genus has no preimage in multiplicity one.
At genus $h$ the remaining regular embedded curve has a single preimage
with trivial automorphism group and complex local degree one. Stability
excludes trees of contracted rational components there. The
compact-support argument on a smooth parameter family representing a
putative relation promotes this local orientation calculation to
independence \cite[Lemmas~5.2--5.3 and Proposition~6.1]{WueProduct}.
Thus we need injectivity of $\Phi$ only on
$V=\bigoplus_h\QQ x_{h,1}$. The argument uses a fixed genus for each
test of a relation.

\begin{lemma}\label{lem:detect}
In $H$, the primitive elements annihilated by every $\rho_l$, $l>1$,
are precisely $V$. Two series with constant coefficient one satisfying
\eqref{eq:localops} through $q^m$ agree to that order if their
images under $\Phi$ do.
\end{lemma}
\begin{proof}
For the first assertion we recall the algebraic argument of
\cite[Lemma~7.1]{WueProduct}. The scalar classes generate $H$ rationally
by \cite[Theorem~6.13]{Pardon}. In the polynomial algebra on formal
scalar generators change variables to the coefficients $v_{h,j}$ of
the logarithm of $1+\sum_{j\ge1}x_{h,j}q^j$. These are primitive, and
$\rho_l v_{h,j}$ is $v_{h,j/l}$ if $l$ divides $j$, and zero otherwise.
These formulas already hold formally. A polynomial relation of smallest
total degree in the new variables has zero constant term. Its reduced
coproduct lies in the tensor square of the space of polynomials of
strictly smaller degree. Evaluation in $H$ is injective on that space,
hence also on its tensor square over $\QQ$. This makes the relation a
primitive polynomial, which is linear in characteristic zero by the
identity $f(v+w)=f(v)+f(w)$. Applying division by its largest label
$j$ would give a nonzero relation among the $x_{h,1}$ (use the identity
if $j=1$). Their smooth images are independent. This proves polynomial
independence of the scalar generators. The same linearity and largest-label
argument in $H$ now identifies its joint primitive kernel with $V$.
For the second assertion, the first coefficient of a difference is
primitive by addition and killed by every proper division. Its smooth
image vanishes, and $\Phi$ is injective on $V$.
\end{proof}

\section{An isolated multiple and the integral basis}\label{sec:basis}

Fix a genus $g$ and a multiplicity $m$. We will compare two
representatives associated with these integers. On the equivariant
side the curve and its balanced bundle are fixed throughout the
construction. On the geometric side the entire holomorphic neighborhood
may depend on $m$.

\begin{lemma}\label{lem:germ}
There is a fixed complex threefold germ about a genus-$g$ curve $C$
with acyclic normal bundle in which every $j[C]$, $1\le j\le m$, is
isolated. Its unit $J_m$ is an integral geometric local curve element of
type $(g,m,0)$, and $J_1=x_{g,1}$.
\end{lemma}
\begin{proof}
For $g\ge2$ take the divisorial deformation of a bundle neighborhood
constructed in the proof of \cite[Lemma~6.7]{Pardon}. Its embedded curve
is unobstructed and has normal bundle of index zero, so it has no
infinitesimal deformations in the fixed threefold. The $m$-multiple lies
in the interior semi-regular locus where the semi-chart of multiples of
deformations of $C$ is open \cite[Lemma~6.3 and Definition~6.6]{Pardon}.
At Chern number zero its fixed-target dimension is zero. After restriction
to a tube retracting onto $C$, every component has Chern number zero.
The point $m[C]$ is therefore open in the entire cycle space of the
germ. The inverse image of that point under addition of $(m-j)[C]$ is
$\{j[C]\}$, proving the lower-degree assertions. The complex point
orientation gives the geometric generator. Deformation to the normal
bundle preserves the regular embedded curve in degree one. Its local
chart is a section over the deformation parameter by the implicit
function theorem. On the normal bundle the unit in degree one is
the class with empty evaluation tuple, namely $x_{g,1}$.

In genus zero take $\cO(-1)^{\oplus2}$. In genus one take
$L\oplus L^{-1}$ with $L$ of infinite order in $\mathrm{Pic}^0(C)$.
A nonzero fiber coordinate on a finite cover is impossible by negative
degree in genus zero. In genus one it would trivialize a pulled-back
line, whose norm would make $L$ torsion. In both cases every compact
curve is supported on the zero section, and $J_j=a_{g,j}=x_{g,j}$ by
the empty-tuple construction.
\end{proof}

\begin{proposition}\label{prop:tube}
For a germ in Lemma~\ref{lem:germ} of genus $g\ge2$, there exist integers
$n_{e,h}$ of finite support, $1\le e\le m$, for which
\begin{equation}\label{eq:tube}
 1+\sum_{j=1}^m J_jq^j =
 \prod_{e,h} A_h(q^e)^{n_{e,h}}\pmod {q^{m+1}},
 \qquad n_{1,h}=\delta_{gh}.
\end{equation}
\end{proposition}
\begin{proof}
We place a small tube of the germ in a closed symplectic six-torus.
In a product of two symplectic elliptic curves take one horizontal torus
and $g-1$ distinct vertical translates. Smoothing the $g-1$ transverse
positive intersections gives a symplectic genus-$g$ surface. Its class
pairs once with a vertical torus, hence stays primitive in the product
with a third elliptic curve. Denote the resulting class by $A$ and take
an integral symplectic form. An orientation-preserving identification
of the surfaces may be used to identify their complex tangent structures
along the core by choosing an appropriate compatible complex structure
on its symplectic tangent plane. Their complex normal bundles are
isomorphic, since both have rank two and degree $2g-2$. Choose a tubular
identification with this normal derivative, taking the complex normal
complement to the symplectic orthogonal. The transported integrable
structure is compatible along the core and tame on a small tube.
Extending a compatible structure along the surface and then using a
cutoff in a chart of tame structures realizes the given germ inside
the closed target.

Set $b=\omega(A)>0$ and bound energy by $mb$. There is a tube $V$
with compact closure in this realization such that the only bounded
nonempty cycles in $\overline V$ are $j[C]$ for $j\le m$. Otherwise
a sequence of tubes shrinking to $C$ contains counterexamples. Retraction
to $C$ and positivity of area give class $jA$ with $1\le j\le m$.
Take a subsequence of fixed class. Bounded-area cycle compactness
\cite[Theorem~2.11]{DIW} forces a limit $j[C]$. Convergence of currents
and supports in the integrable tube contradicts isolation in the
holomorphic cycle space. Inside such an isolating tube take $V$ with
compact closure and with $C$ in its interior. Every bounded cycle
contained in $\overline V$ lies in the larger isolating tube, so none
can meet $\partial V$. This property persists under small
perturbations by the same compactness argument. Containment in $V$
therefore defines a compact open and closed locus over a short closed
path to a generic tame structure. Its unit transports $\Phi(J_j)$
in each class $jA$.

At the generic endpoint, the bounded connected supports are disjoint
embedded curves. Partition the supports in $V$ into finitely many
clusters, including the exterior as a separate part
\cite[Propositions~3.18 and~3.21]{DIW}. In each part the union of
the underlying curves is compact by closed incidence. Distinct parts have
disjoint unions, and hence disjoint open target neighborhoods. A
disconnected cycle then decomposes uniquely into cycles in the respective
neighborhoods, including empty cycles. The unit for the union is the
product of the corresponding series, with the convolution performed in
integral homology. This is the excision argument of \cite[\S4]{WueProduct}.
When a path is considered in the ambient threefold, we keep only the
finitely many labels occurring below the energy cutoff. Indeed the
bounded cycle space over the closed path is compact and the integral
homology label is locally constant. Products and inverses of these
series still have finite support at the cutoff, since their positive
degrees have a positive minimum.

For one of these clusters, of class $eA$, apply the cluster isotopy theorem
\cite[Theorem~3.25]{DIW} in the geometric form of
\cite[\S4, proof of Proposition~4.1]{WueProduct}. On a regular arc a
subdivision changes the chosen factor only by factors in proper
multiples. At a fold, compact continuation of the outer insertion gives
the product of the two inner factors equal to one modulo such corrections.
Consequently a crossing inverts the distinguished factor, modulo a
finite product of corrections. At the elementary endpoint its factor is
$\Phi(A_h(q^e))$. Iterate on the correction clusters; their core classes
are $keA$, $k>1$, and $\lfloor m/(ke)\rfloor<\lfloor m/e\rfloor$.
This gives a finite signed product in $S$ at the required cutoff. The
use of disjoint target neighborhoods in excision is what allows these
geometric identities to include the cycles with several components.

Addition and division for the isolated units satisfy
\eqref{eq:localops} up to $m$: the decompositions of $j[C]$ are precisely
$i[C]+(j-i)[C]$. The product on the right of \eqref{eq:tube} obeys
the same laws because substitution, multiplication and inversion
preserve them. Lemma~\ref{lem:detect} lifts the identity from $S$ to $H$.
In degree one it gives $x_{g,1}=\sum_h n_{1,h}x_{h,1}$, so independence
makes the asserted leading exponent positive and equal to one.
\end{proof}

\begin{proof}[Proof of Theorem~\ref{thm:basis}]
Take the coefficient of $q^m$ in \eqref{eq:tube}. The unique unshifted
factor is $A_g(q)$. Any other factor starts in degree $e\ge2$ and only
uses coefficients of multiplicities at most $\lfloor m/e\rfloor<m$.
The binomial coefficients in all integer powers, including negative
powers, are integers. Hence $J_m=a_{g,m}+P_{g,m}$ with $P_{g,m}$
an integral polynomial in lower multiplicities. Choose one germ for
each pair $(g,m)$ as in Lemma~\ref{lem:germ}. Pardon's integral
generation theorem permits \emph{any one such choice for each type}
\cite[Theorem~6.13]{Pardon}. Nonnegative Chern numbers exclude other
types from this summand. Induction in $m$, simultaneously for all $g$,
identifies the subalgebra of $H$ generated by these integral images with
$\ZZ[a_{g,m}]$.

To see that it is a polynomial ring, apply the least-degree relation
argument of Lemma~\ref{lem:detect} to formal balanced generators.
Their logarithmic coordinates are primitive and have the same division
law, and their degree-one images are again the independent $x_{g,1}$.
There is no rational polynomial relation. The triangular equations with
leading coefficient one imply that the chosen $J_m$ for all types
also have algebraically independent images over $\QQ$. Since those
classes generate the original integral summand, its rationalization map
has zero kernel. This proves the stated integral identification.
\end{proof}

\section{Integral products on compact threefolds}\label{sec:compact}

Let $\Gamma$ be an additive commutative monoid with an additive degree
$E:\Gamma\to\mathbb R_{\ge0}$, positive away from zero, with finite
sublevel sets. Multiplication in $H[[\Gamma]]$ is the convolution for
addition in $\Gamma$. An element has only finitely many decompositions
into nonzero summands: at its degree there are finitely many possible
sublabels and these have a positive minimum degree. For a series
$P=\sum P_\beta Q^\beta$ with $P_0=1$, compatibility with the
two geometric operations means
\begin{equation}\label{eq:twolaws}
 \Delta P_\beta=\sum_{\gamma+\delta=\beta}P_\gamma\otimes P_\delta,
 \qquad
 \rho_l P_\beta=\sum_{l\gamma=\beta}P_\gamma\quad(l>1).
\end{equation}
The second sum ranges over every root in $\Gamma$. Distinct integral
classes may have the same multiple.

\begin{proposition}\label{prop:criterion}
A compatible series has all its coefficients in $H_{\ZZ}$ if and only if it has
a product expansion in the factors $A_h(Q^\beta)$ with integral
exponents, finitely supported in $h$ at each $\beta\ne0$. In that case
the exponents are unique.
\end{proposition}
\begin{proof}
A single factor satisfies \eqref{eq:twolaws} by \eqref{eq:localops}.
Equivalently the second identity says that applying $\rho_l$
coefficientwise gives the substitution $Q^\gamma\mapsto Q^{l\gamma}$;
this identity and the coproduct identity persist under products and
inversion. All factors and their inverses have integral coefficients
by Theorem~\ref{thm:basis}.

Conversely divide out the factors already found at smaller degree.
At a first remaining degree the coefficient $R_\beta$ of the residual
integral series is primitive. Indeed in the first sum of
\eqref{eq:twolaws} every proper summand has smaller degree. For the
same reason every term on the right of the second identity vanishes.
Lemma~\ref{lem:detect} puts $R_\beta$ in $V$. By the integral
polynomial basis and $a_{h,1}=x_{h,1}$ we have
$V\cap H_{\ZZ}=\bigoplus_h\ZZ x_{h,1}$. The coefficients in that direct
sum are precisely the required exponents at $\beta$. Dividing out the
corresponding factors changes no other label of this degree.
Finite sublevel sets give the expansion and the same recursion proves
uniqueness.
\end{proof}

\begin{proof}[Proof of Theorem~\ref{thm:compact}]
On a compact Kähler threefold the mass of an effective cycle for the
Kähler metric is its area. Bounded-area cycles form a compact space
and their integral homology label is locally constant in the geometric
topology \cite[Theorem~2.11 and Remark~2.13]{DIW}. The image of a
compact space in this discrete group is finite, even when the group has
torsion. The cycles of fixed class give a compact open and closed
locus. A nonzero effective class has positive area. Nonnegativity of
$c_1(TX)$ implies that all components of a cycle in $\Gamma_0(X)$
have zero first Chern number. In particular the unit of the whole
locus is an integral semi-Fano insertion.

Addition of cycles pulls this unit back to the sum of the exterior
products of the units on the class loci for all splittings of $\beta$.
Multiplication by $l$ pulls it back to the units on all class loci
with $l\gamma=\beta$. These are precisely \eqref{eq:twolaws}, so
Proposition~\ref{prop:criterion} proves \eqref{eq:compact}. On a compact
manifold the areas of complex tangent lines under two Kähler metrics
are uniformly comparable. This identifies the two completions, and
uniqueness gives the same exponents. Under the projective hypotheses,
the geometric product of \cite[Theorem~1.1]{WueProduct} has exactly these
local factors and signed cluster multiplicities. Its rationalization
agrees with \eqref{eq:compact}; uniqueness identifies the integers.
\end{proof}

\section{Rational curves on a contracted face}\label{sec:ruling}

For a genus-zero curve the scalar and balanced factors coincide. Write
$A_0(q)=\sum_{m\ge0}x_{0,m}q^m$. A family of resolved nodes shows how
a negative power of this factor can arise even though every cycle is
effective. We give the family first because it also computes the
contribution of a moving exceptional ruling.

Let $E$ be the nonsplit extension of $\cO$ by $\cO(-2)$ over $\PP^1$.
Then $E\simeq\cO(-1)^{\oplus2}$, since the boundary map from
$H^0(\cO)$ to $H^1(\cO(-2))$ is an isomorphism. Let
$u:\tot E\to\CC$ be the quotient coordinate. For a smooth complex
function $f$ on a closed disc, nonzero on the boundary, form the
fiber product $u=f(t)$ over its interior. A complex-horizontal
right inverse $h$ to $du$ gives this threefold an almost-complex
structure: in the splitting of a tangent vector as
$(a,k+h(df(a)))$ use multiplication by $i$ on $a$ and the vertical
complex structure on $k$. The projection to the disc is holomorphic
and each surface fiber retains its original complex structure.
For holomorphic $f$ this is the complex fiber product.

\begin{lemma}\label{lem:fold}
The series of compact cycles in this family, graded by their degree
over $\PP^1$, is $\Phi(A_0(q))^{\deg(f,0)}$ in the smooth algebra.
For holomorphic $f$ it is $A_0(q)^{\deg(f,0)}$ in $H$.
\end{lemma}
\begin{proof}
A compact irreducible curve projects constantly to the disc. The only
compact irreducible curve in $\tot E$ is its zero section, by negativity
after pullback to the normalization. Thus over a zero of $f$ there is
one rational curve, and there are no others. The full degree-$m$ cycle
locus is $\Sym^m(f^{-1}(0))$: the parametrization by sums is continuous
in support and current and is a continuous bijection from a compact
space to a Hausdorff space. The same construction over a compact path
without boundary zeros transports its unit. For separated zero sets
one can excise into disjoint targets, so their series multiply.

A positive simple zero can be deformed to the identity map, and there
the family is an open neighborhood of the zero section of $E$.
Its factor is $\Phi(A_0)$. At a negative regular zero the germ deforms
to complex conjugation by linearizing at the zero and moving within
the orientation-reversing linear isomorphisms. The local homotopies
avoid boundary zeros on a small disc. In the family
$f_s(x+iy)=x^2-s+iy$ there are no zeros for $s<0$ and two regular
zeros of opposite signs for $s>0$. The two factors therefore multiply
to one. Translation of any $f$ by a sufficiently small regular value
counts its simple zeros with Brouwer signs, proving the assertion.
For holomorphic $f$ the small translations are a holomorphic family
and the simple zeros are positive; the same computation is made
directly on complex cycles.
\end{proof}

For $f(t)=t^k-s$, continuation around $s=0$ permutes the $k$ supports
cyclically. The limiting single support has factor $A_0^k$. In bundle
coordinates its transition at $s=0$ is
$v'=z^2v+zt^k$, $t'=t$. Thus for distinct $k,l>r$ the $r$th
infinitesimal neighborhoods of the zero curves agree while their series
differ already in degree one. This is a dependence on the nonlinear
neighborhood that persists however large the prescribed jet order is.

\begin{proposition}\label{prop:ruled}
Let $p:X\to W$ be a birational contraction from a smooth projective
threefold with $K_X\simeq\cO_X$ to a projective variety. Suppose $p$ is an
isomorphism off a smooth divisor $D=\PP_B(V)$, and contracts exactly
its ruling fibers to the smooth connected curve $B\simeq p(D)$. If $b=g(B)$ and
$\ell$ is a fiber class, the complete universal series on the contracted
face is $A_0(Q^\ell)^{2b-2}$.
\end{proposition}
\begin{proof}
The pullback of an ample divisor on $W$ exposes the contracted face.
Every irreducible component there is a ruling fiber, and in class
$m\ell$ the cycle locus is $\Sym^m B$. All fibers have the same
integral class because $B$ is connected. An ample degree on $X$
distinguishes its multiples, even when $\ell$ is divisible in
$H_2(X;\ZZ)$. Deform $D$ to its normal
bundle in $X$ by blowing up $D\times0$ in $X\times\CC$ and deleting
the strict transform of the central fiber. The normal bundle is $K_D$.
A compact curve on the special fiber with contracted pushforward lies
over one point of $B$. On that ruling the normal bundle is $\cO(-2)$,
so the curve is the zero section. The relative fixed-label locus over
a closed parameter interval remains $\Sym^m B$ times that interval,
by the same compact continuous parametrization as in Lemma~\ref{lem:fold}.

Use the convention of lines in $V$. The relative Euler sequence gives
\[
 0\longrightarrow K_D\longrightarrow
 \cO_D(-1)\otimes\pi^*(V^\vee\otimes K_B)
 \xrightarrow{u}\pi^*K_B\longrightarrow0 ,
 \qquad \pi:D\to B.
\]
Replace the zero section on the right by a transverse smooth section
of $K_B$. Locally this is the family in Lemma~\ref{lem:fold}, since
on each ruling the Euler extension is nonsplit. A complex-linear
horizontal splitting gives the same fiberwise structures as in that
lemma and can be exchanged with a local product splitting by a homotopy.
Globally the projection to $B$ remains pseudoholomorphic. Each curve
of zero base degree therefore stays in one fiber; along the homotopy
of sections the relative symmetric powers of the compact zero loci
give compact support in every degree. The signed number of zeros of
the section is $2b-2$. Multiplying the local factors proves the assertion
after applying $\Phi$. Both complex series satisfy addition and
division in the integer label $m$, since every effective summand or
root of a contracted cycle is contracted. Lemma~\ref{lem:detect}
lifts the equality. For $b\ge1$ one can equally use a nonzero
holomorphic differential and count its zeros holomorphically.
The cotangent obstruction and its Euler number in the numerical
type III theory are due to Wilson \cite[\S4]{Wilson}.
\end{proof}

\section{Specialization across a conifold transition}\label{sec:transition}

Suppose a flat projective family over a smooth curve has smooth
algebraic total space of dimension four and central fiber with only
ordinary double points. Let $X_s$ be a smooth nearby fiber and let
$Y\to X_0$ be a projective small resolution. Assume the two smooth
threefolds have trivial canonical bundles and torsion-free second
integral homology. Fix the common-complement marking and a relatively
ample line bundle. Write $\pi:H_2(Y;\ZZ)\to H_2(X_s;\ZZ)$ for the
resulting quotient by the exceptional curve classes $[C_i]$.
Writing $U_X$ for the series in \eqref{eq:compact}, the universal
specialization identity of \cite[Theorem~3.2]{WueConifold} reads
\begin{equation}\label{eq:transition}
 U_{X_s}=\pi_*\left(U_Y\prod_i A_0(Q^{[C_i]})^{-1}\right).
\end{equation}
Here $\pi_*$ adds coefficients of classes with the same image.
Finiteness requires the exceptional normalization inside the parentheses.
In the blowup $\phi:\widetilde Y\to Y$ along the $C_i$, put
$j=\mathrm{PD}_{\widetilde Y}^{-1}\phi^*\mathrm{PD}_Y$. After normalization
all nonzero labels belong to
$F=\{\gamma\ne0:j\gamma\text{ is effective on }\widetilde Y\}$.
This set is closed under addition, avoids $\ker\pi$, and has finite
sublevel sets for the projected ample degree
\cite[Proposition~3.1]{WueConifold}. Indeed the ample degree of $j\gamma$
for a suitable polarization on the blowup is a fixed positive
multiple of the projected degree, since $j\gamma$ pairs to zero with
every exceptional divisor. The support assertion in \eqref{eq:transition}
in that generality uses the numerical reconstruction of
\cite[\S2]{WueConifold}. It assumes that the indicated connected unmarked
Gromov--Witten evaluations descend to the relevant relations in the
rational cycle algebra, including those arising from families of
complex targets. The analytic-family input to that descent is discussed
in \cite[Remark~1.8]{Pardon}. There is a direct argument using only
fixed-target counts when both endpoints also have
$H^1(\cO)=H^2(\cO)=0$.

The normalized series has integral balanced exponents
$\widetilde n_{\gamma,h}$, obtained by subtracting one at each
exceptional factor with its actual class. On the latter projective
domain let $G_h(q,u)$ denote the connected numerical series of the
elementary genus-$h$ curve. The signed-cluster comparison of
\cite[\S7.4]{WueProduct} identifies the connected invariants of each fixed
endpoint with the sum of $n_{\gamma,h}G_h(Q^\gamma,u)$. Write
$w=(2\sin(u/2))^2$; the coefficient of $q$ in $G_h$ is $w^{h-1}$.
The contracted face consists exactly of sums of the disjoint resolved
nodes, so subtracting their $G_0$ removes all connected contributions
on that face. For $\pi\gamma\ne0$ the ordinary degeneration of the
resolution gives its connected invariant as the empty-contact
relative invariant of $j\gamma$ on the blowup. The nonempty contact
strata have negative virtual dimension on the blowup side. In
particular every connected coefficient at $\gamma\notin F$ vanishes
after normalization \cite[\S2.1]{LLW}, \cite[\S3]{WueConifold}.

Induct on an ample degree on the resolution. If $\gamma\notin F$,
every proper root of $\gamma$ is outside $F$, since $F$ is closed
under multiplication by positive integers. After the lower-degree
terms have vanished, the connected coefficient at $\gamma$ is
$\sum_h\widetilde n_{\gamma,h}w^{h-1}=0$. This is a finite Laurent
polynomial and its least power has a distinct leading order in $u$.
All its coefficients therefore vanish. Thus every remaining exponent
is supported in $F$, and its product may be pushed forward in finite
sums. The connected smoothing formula in the total-space class adds
the fixed-target resolution counts with that image, also for classes
ineffective on the smoothing \cite[\S3, proof of Theorem~3.2]{WueConifold},
\cite[Theorem~2.16 and \S3]{LW}. Comparing with the pushed-forward
product, the same induction now in the smoothing degree identifies
its exponents. This proves \eqref{eq:transition} on the stated
projective domain without a theory of evaluations on general families.
With the numerical reconstruction assumed in the broader case, the
support theorem in \cite[\S3]{WueConifold} and successive division by
factors in $F$ give the same exponent formula. In either case,
for every nonzero effective $\beta$ on the smoothing,
\[
 n_{\beta,h}(X_s)=\sum_{\substack{\gamma\in F\\\pi\gamma=\beta}}
                       \widetilde n_{\gamma,h}.
\]
The proof applies in the monoid generated by the effective smoothing
classes and $\pi(F)$, extending the series by zero. This also describes
the cancellation when an image class is ineffective. Two small
resolutions have the same blowup and their normalized series agree
under its common zero-contact identification \cite[\S3]{WueConifold}.
The same summation for scalar product exponents preserves their
integrality whenever the source exponents are integral
\cite[Corollary~3.3]{WueConifold}. The scalar exponent in
\S\ref{sec:intro} explains the need for this hypothesis.

One can bound the sum explicitly on the classical transition of a
quintic containing a plane \cite{Rossi}.
Let $f,g$ be general quartics on $\PP^4$. The central quintic
$x_3f+x_4g=0$ has sixteen nodes on the plane $x_3=x_4=0$.
Indeed the two restricted quartics meet transversely, and at each
intersection $x_3,x_4,f,g$ are local coordinates. Away from the
plane smoothness follows from Bertini. In $\PP^1_{[s:t]}\times\PP^4$
the equations $s x_3+t x_4=0$ and $-s g+t f=0$ give its projective
small resolution $Y$. The projection selects the kernel line of the
two-by-two coefficient matrix, with an entire ruling over each node.
Adjunction gives $K_Y=\cO_Y$. Adding a general multiple of a quintic
nonvanishing at all sixteen nodes gives a smoothing with smooth total
space near the central fiber. Write $B,H$ for the hyperplane classes
from the two factors. Integral Lefschetz for the two ample divisors gives
coordinates $(a,d)=(B\cdot\gamma,H\cdot\gamma)$ on $H_2(Y;\ZZ)$,
and the smoothing map is $(a,d)\mapsto d$. Each exceptional curve
has class $(1,0)$.

In this example $F$ satisfies $0\le a\le 4d$, $d\ge1$. To see the
upper bound, write $u=f/s=g/t$ and $v=x_4/s=-x_3/t$. These sections
have classes $4H-B$ and $H-B$. The sections $u,vh$, with $h$ ranging
over cubics, generate the ideal of the exceptional curves after twisting
by $4H-B$. Indeed their common ideal is locally $(u,v)$, the two normal
coordinates in the small resolution of the node. Thus on the blowup
$\phi^*(4H-B)-E$ is generated by global sections, where $E$ is the
sum of the exceptional divisors. Pairing this nef divisor and $\phi^*B$
with $j\gamma$ proves both inequalities. A line in the plane missing
the nodes lifts to the graph of the degree-four map $[f:g]$ on the
line. Its proper transform has class $j(4,1)$, so the upper bound is
attained in $F$. After division by $A_0(q)^{16}$ the coefficient in
smoothing degree $d>0$ is therefore a polynomial in the exceptional
variable $q$ of degree at most $4d$. Membership in this support bound
alone does not establish that a coefficient is nonzero.

In the four-node toric transition let $\beta=(1,1,0)$ be the
primitive line class of the smoothing plane. In the resolution basis
of \cite[\S4]{WueConifold} the marking takes
$(n_1,\ldots,n_6)$ to $(n_2,n_3,n_4)$. The following individual source and smoothing coefficients use
the stable-pair evaluation of geometric local classes assumed in
\cite[\S\S5--6]{WueConifold}; see also \cite[Remark~1.8]{Pardon}.
The conifold specialization proved above uses the fixed-target
Gromov--Witten comparison. Of the 37 candidate lifts of
$\beta$, eight contribute. Their complete coefficients are given in
\cite[(21)--(24)]{WueConifold}. The reconstruction also uses a bound on the Euler characteristic of a pure curve. The quadratic terms cancel upon specialization, so
the eight polynomials sum to $3x_{0,1}$. On the smoothing the minimal
pair spaces for lines and conics are the dual plane and $\PP^5$.
Their signed Euler characteristics $3$ and $-6$ give
\[
 U_{X_s,\beta}=3x_{0,1},\qquad
 U_{X_s,2\beta}=3x_{0,2}+3x_{0,1}^2-6x_{0,1}
 \qquad\text{\cite[Proposition~6.1]{WueConifold}}.
\]
The line spans an extremal ray and the only effective splitting of its
double is $\beta+\beta$. Thus the first factor on the ray is
$A_0(Q^\beta)^3$. At $2\beta$ its contribution is
$3x_{0,2}+3x_{0,1}^2$. The remaining $-6x_{0,1}$ is the leading
coefficient of $A_0(Q^{2\beta})^{-6}$.

\begin{thebibliography}{Wue26e}
\bibitem[BP08]{BP} J.~Bryan and R.~Pandharipande,
\emph{The local Gromov--Witten theory of curves}, J. Amer. Math. Soc.
\textbf{21} (2008), 101--136,
\href{https://doi.org/10.1090/S0894-0347-06-00545-5}{doi:10.1090/S0894-0347-06-00545-5}.
\bibitem[DIW26]{DIW} A.~Doan, E.-N.~Ionel and T.~Walpuski,
\emph{The Gopakumar--Vafa finiteness conjecture}, Ann. of Math. (2)
\textbf{204} (2026), 819--871,
\href{https://doi.org/10.4007/annals.2026.204.3.1}{doi:10.4007/annals.2026.204.3.1}.
Numbering follows \href{https://arxiv.org/abs/2103.08221v2}{arXiv:2103.08221v2}.
\bibitem[IP18]{IP} E.-N.~Ionel and T.~H.~Parker,
\emph{The Gopakumar--Vafa formula for symplectic manifolds},
Ann. of Math. (2) \textbf{187} (2018), 1--64,
\href{https://doi.org/10.4007/annals.2018.187.1.1}{doi:10.4007/annals.2018.187.1.1}.
\bibitem[LLW18]{LLW} Y.-P.~Lee, H.-W.~Lin and C.-L.~Wang,
\emph{Towards $A+B$ theory in conifold transitions for Calabi--Yau threefolds},
J. Differential Geom. \textbf{110} (2018), 495--541,
\href{https://arxiv.org/abs/1502.03277v3}{arXiv:1502.03277v3}.
\bibitem[LW25]{LW} Y.~Lin and S.-S.~Wang,
\emph{Gromov--Witten/Pandharipande--Thomas correspondence via conifold transitions},
Trans. Amer. Math. Soc. \textbf{378} (2025), 4413--4436,
\href{https://arxiv.org/abs/2310.18170v3}{arXiv:2310.18170v3}.
\bibitem[Par25]{Pardon} J.~Pardon, \emph{Universally counting curves in
Calabi--Yau threefolds}, \href{https://arxiv.org/abs/2308.02948v3}{arXiv:2308.02948v3} (2025).
\bibitem[Ros06]{Rossi} M.~Rossi, \emph{Geometric transitions},
J. Geom. Phys. \textbf{56} (2006), 1940--1983,
\href{https://doi.org/10.1016/j.geomphys.2005.09.005}{doi:10.1016/\allowbreak j.geomphys.2005.09.005}.
\bibitem[Wil99]{Wilson} P.~M.~H.~Wilson,
\emph{Flops, type III contractions and Gromov--Witten invariants on
Calabi--Yau threefolds}, in \emph{New Trends in Algebraic Geometry},
London Math. Soc. Lecture Note Ser. \textbf{264} (1999), 465--484,
\href{https://doi.org/10.1017/CBO9780511721540.018}{doi:10.1017/CBO9780511721540.018}.
\bibitem[Wue26a]{WueCounter} B.~J.~Wuebben,
\emph{A counterexample to integral product exponents in universal curve counting}, 28 October 2026, revised 21 November 2026,
\href{https://doi.org/10.5281/zenodo.23750048}{doi:10.5281/zenodo.23750048}.

\bibitem[Wue26b]{WueProduct} B.~J.~Wuebben,
\emph{A geometric product formula for universal curve counting on Calabi--Yau threefolds}, 23 November 2026,
\href{https://doi.org/10.5281/zenodo.24455518}{doi:10.5281/zenodo.24455518}.

\bibitem[Wue26c]{WueBalanced} B.~J.~Wuebben,
\emph{Antidiagonal local curves and integral exponents in universal curve counting}, 11 November 2026, revised 22 November 2026,
\href{https://doi.org/10.5281/zenodo.24432700}{doi:10.5281/zenodo.24432700}.

\bibitem[Wue26h]{WueTriples} B.~J.~Wuebben,
\emph{Triple curves and unbounded denominators in threefold curve counting}, 24 November 2026,
\href{https://doi.org/10.5281/zenodo.24510966}{doi:10.5281/zenodo.24510966}.

\bibitem[Wue26i]{WueConifold} B.~J.~Wuebben,
\emph{Universal curve counting under conifold transitions}, 28 October 2026, revised 21 November 2026,
\href{https://doi.org/10.5281/zenodo.23750046}{doi:10.5281/zenodo.23750046}.

\end{thebibliography}

\end{document}
